\scanpage{48}
% archive foliation: 47
3) Relations entre les deux théories (pour schémas abéliens et pour \deleted{\unsure{2)}Relations} groupes de BT). Supposons que p soit localement nilpotent sur S. Travaillons \deleted{On suppose} à nouveau avec le site cristallin de \unsure{Berthel\deleted{l}ot} (pas le nilpotent). Le résultat essentiel sera, pour un schéma abélien variable\deleted{\unsure{r}} A sur S, un isomorphisme fonctoriel compatible aux changements de base
\[
\underline{D}^{\inserted{*}}(A(\infty))\simeq R^{1}f_{\mathrm{cris}\inserted{*}}(\underline{O}_{A_{\mathrm{cris}}}),
\]
cet isomorphisme induisant, pour les valeurs des deux membres sur S, un isomorphisme compatible aux filtrations et (compte tenu des isomorphismes $\underline{t}_{A(\infty)}\simeq\underline{t}_A$ et idoine pour $A^{*}$).\inserted{-- C'est en fait en postulant un tel isomorphisme qu'on arrive par voie heuristique à une définition de $\mathbb{D}^{*}(G)$ pour un groupe de BT G quelconque.}

Ceci impliquera en particulier le fait que, pour un \deleted{voisinage} épaississement à p. div. loc. nilp. \deleted{de} S' de S, il y a identité de la théorie des prolongements infinitésimaux de A, et de $A(\infty)$. De \unsure{eeci|ceci}, par dévissage sur l'ordre de nilpotence, on peut déduire par exemple immédiatement le théorème de Serre-Tate déjà annoncé plus haut:\inserted{méthode suivie par Messing ds sa thèse.}

[On peut aussi définir une extension vectorielle universelle $\deleted{\unsure{\underline{E}}}\ E(G)$ d'un groupe de BT, et on trouvera une isomorphisme canonique etc
\[
E(A(\infty))\simeq E(A)(\infty)
\]
donnant naissance, par le procédé habituel, à un isomorphisme plus général de cristaux en groupes
\[
\underline{E}(A(\infty))\simeq\underline{E}(A)(\infty)\quad ....
\]
]

\deleted{4) Passage à la limite}

\deleted{4) Problème: trouver le foncteur des F-cristaux filtrés vers \deleted{\unsure{modules}} faisceaux p-adiques.}

